\documentclass{ip-lab}

\lablevel{n2}
\labnumber{Laboratorio 3}
\labtitle{Diccionarios y cadenas de caracteres}

\makeheader

\begin{document}

\maketitle

\begin{sectionbox}{Preparación del ambiente de trabajo}
  Cree un nuevo módulo en el que escribirá las funciones correspondientes a las actividades 1 a 3 de este laboratorio.
\end{sectionbox}

\begin{sectionbox}{Actividad 1: Palíndromos}
Escriba una función que reciba cuatro números de exactamente 6 dígitos y determine cuáles son palíndromos. Un número es palíndromo si se lee igual de izquierda a derecha que de derecha a izquierda.

La función debe retornar un diccionario donde cada número sea una llave y su valor sea un booleano indicando si es palíndromo.

\textbf{E.g.:} Para los números \texttt{123321}, \texttt{988789}, \texttt{938726} y \texttt{927729} la función debería retornar:

\texttt{\{123321: True, 988789: False, 938726: False, 927729: True\}}
\end{sectionbox}

\begin{sectionbox}{Actividad 2: Competencia de videojuegos}
En un torneo de videojuegos, cuatro jugadores compiten por ser el campeón. Cada jugador tiene información almacenada en un diccionario con las siguientes llaves: \texttt{\textquotedbl nickname\textquotedbl} (nombre de usuario, sin espacios ni tildes), \texttt{\textquotedbl puntaje\textquotedbl} (puntaje obtenido), \texttt{\textquotedbl experiencia\textquotedbl} (años de experiencia), \texttt{\textquotedbl plataforma\textquotedbl} (consola utilizada), y \texttt{\textquotedbl region\textquotedbl} (región geográfica).

Implemente una función que reciba cuatro jugadores como parámetros y retorne el nickname del jugador con el puntaje más alto. En caso de empate en el puntaje, priorice al jugador que tenga más consonantes en su nickname y, si persiste el empate, priorice al jugador que tenga menos experiencia. Se garantiza que no habrá empates completos.

\textbf{Nota:} Las vocales en español son: a, e, i, o, u; el resto de letras son consonantes.
\end{sectionbox}

\pagebreak

\begin{sectionbox}{Actividad 3: Sistema de códigos de productos}
Una tienda en línea necesita validar códigos de productos que deben cumplir las siguientes reglas:
\begin{itemize}
  \item Deben tener exactamente 4 caracteres de longitud.
  \item No deben tener letras mayúsculas.
  \item Los dígitos del código deben estar ordenados estrictamente ascendentemente.
\end{itemize}

Implemente una función que reciba cuatro códigos de productos como parámetros y retorne un diccionario donde cada código sea una llave y su valor sea un booleano indicando si el código es válido.

\textbf{E.g.:} Para los códigos \texttt{\textquotedbl 1a23\textquotedbl}, \texttt{\textquotedbl x0y\textquotedbl}, \texttt{\textquotedbl 4B56\textquotedbl}, \texttt{\textquotedbl 8m79\textquotedbl}, la función debería retornar:

\texttt{\{\textquotedbl 1a23\textquotedbl: True, \textquotedbl x0y\textquotedbl: False, \textquotedbl 4B56\textquotedbl: False, \textquotedbl 8m79\textquotedbl: False\}}
\end{sectionbox}

\begin{sectionbox}{Entrega}
Entregue el archivo a través de Bloque Neón en el laboratorio del Nivel 2 designado como ``N2-L3: Diccionarios y cadenas de caracteres''.
\end{sectionbox}

\end{document}
